\chapter{Формализация задачи обеспечения согласованности игровых данных}

\section{Математическое представление игрового состояния}

Игровой мир многопользовательского приложения можно представить в виде динамической системы,
состояние которой изменяется во времени под воздействием действий игроков и игровой логики. 
Состояние мира в момент времени $t$ обозначим как $S(t)$:

\[
S(t) = \{ s_1(t), s_2(t), \ldots, s_n(t) \},
\]

где $s_i(t)$ - состояние $i$-ой игровой сущности (позиция, скорость, параметры и другие).

Каждое состояние можно формально представить как вектор параметров:

\[
s_i(t) = \left( p_i(t), v_i(t), a_i(t), \lambda_i(t) \right),
\]

где:
\begin{itemize}
    \item $p_i(t)$ - позиция сущности,
    \item $v_i(t)$ - скорость,
    \item $\lambda_i(t)$ - дискретные состояния (статусы, флаги).
\end{itemize}

Все состояние мира представляет собой многомерное пространство, которое необходимо корректно поддерживать на клиенте и сервере.

\section{Модель клиент-серверного взаимодействия}

В большинстве сетевых игр сервер является единственным источником истинного состояния. 
Клиентская сторона может оперировать временными предсказаниями, но окончательные значения 
всегда принадлежат серверу.

На стороне клиента поддерживается приближенное состояние $\hat{S}(t)$, которое может отличаться 
от серверного $S(t)$. Клиент получает обновления состояния в дискретные моменты времени $t_k$:

\[
S(t_k) \rightarrow \text{Клиент},
\]

а отправляет данные о действиях игрока в форме входных команд:

\[
u(t_k) = \text{input}(t_k), \qquad u(t) \in U,
\]

где $U$ - пространство возможных действий.

Сервер выполняет симуляцию:

\[
S(t + \Delta t) = F(S(t), u_1(t), u_2(t), \ldots, u_m(t)),
\]

где $F$ - функция игровой логики, $m$ - количество игроков.

\section{Определение рассогласования состояния}

Пусть клиентское состояние $\hat{S}(t)$ является приближением серверного состояния $S(t)$.
Тогда показатель рассогласования:

\[
D(S(t), \hat{S}(t)) = \sum_{i=1}^{n} w_i \cdot d(s_i(t), \hat{s}_i(t)),
\]

где:
\begin{itemize}
    \item $d()$ - расстояние между состояниями одной сущности,
    \item $w_i$ - вес значимости сущности.
\end{itemize}

$d()$ включает позиционную ошибку:

\[
d(s_i, \hat{s}_i) = 
    \alpha \cdot \|p_i - \hat{p}_i\| \;+\;
    \beta \cdot \|v_i - \hat{v}_i\|\;+\;
    \gamma \cdot \delta(\lambda_i, \hat{\lambda}_i),
\]

где:
\begin{itemize}
    \item $\|p_i - \hat{p}_i\|$ - расстояние между фактической позицией сущности 
    и её клиентским приближением; характеризует пространственную ошибку;
    \item $\|v_i - \hat{v}_i\|$ - разница скоростей сущности, отражающая рассогласование 
    динамики движения;
    \item $\delta(\lambda_i, \hat{\lambda}_i)$ - бинарная функция различия дискретных параметров:
    \[
        \delta(\lambda_i, \hat{\lambda}_i) =
        \begin{cases}
            0, & \text{если } \lambda_i = \hat{\lambda}_i,\\
            1, & \text{если } \lambda_i \neq \hat{\lambda}_i;
        \end{cases}
    \]
    \item $\alpha$, $\beta$, $\gamma$ - вес значимости соответствующих ошибок.
\end{itemize}


Допустимый порог рассогласования определим как:

\[
D(S(t), \hat{S}(t)) \leqslant D_{\text{max}}.
\]

Если ошибка выше порога, то требуется корректировка состояния.

\section{Постановка задачи обеспечения согласованности}

Целью системы синхронизации является минимизация рассогласования между клиентом и сервером при ограничениях на сетевые ресурсы и задержки:

\[
\min_{\Pi} \; D\left(S(t), \hat{S}_{\Pi}(t)\right),
\]

где $\Pi$ - выбранный алгоритм синхронизации.

При этом должны выполняться следующие ограничения:

\[
B(\Pi) \leqslant B_{\text{max}}, \qquad 
L(\Pi) \leqslant L_{\text{max}},
\]

где:
\begin{itemize}
    \item $B(\Pi)$ - требуемая пропускная способность,
    \item $L(\Pi)$ - задержка реакции,
    \item $B_{\text{max}}$ и $L_{\text{max}}$ - допустимые ограничения сети.
\end{itemize}

В общем виде задача синхронизации представляет собой многокритериальную оптимизацию:

\[
\min_{\Pi} \left(
    \alpha \cdot D +
    \beta \cdot B +
    \gamma \cdot L
\right),
\]

где $\alpha, \beta, \gamma$ - веса важности согласованности, трафика и задержки.

\section{Требования к алгоритмам синхронизации}

Алгоритмы, обеспечивающие согласованность, должны удовлетворять следующим требованиям:

\begin{itemize}
    \item Точность - алгоритм должен обеспечивать минимальные расхождения 
    между состоянием на клиенте и сервере, чтобы игроки наблюдали корректную картину происходящего.
    \item Стабильность - изменения состояния и корректировки со стороны сервера 
    должны происходить плавно, без резких скачков и визуальных искажений.
    \item Устойчивость к сетевым ошибкам - система должна корректно работать при задержках, потере отдельных пакетов и ухудшении качества соединения.
    \item Эффективность - алгоритм не должен перегружать вычислительные ресурсы 
    и должен использовать как можно меньше сетевого трафика.
\end{itemize}

Задача обеспечения согласованности игровых данных формализуется как задача оптимизации, уравновешивающая точность, скорость обновления и сетевые ограничения.